\documentclass[11pt]{article}
\usepackage[margin=1in]{geometry}
\usepackage{amsmath,amssymb,amsthm,booktabs,longtable}
\usepackage[hidelinks]{hyperref}
\newtheorem{lemma}{Lemma}
\newcommand{\ii}{\mathrm i}
\newcommand{\ZS}{\mathcal S}
\newcommand{\PS}{\mathcal P}
\newcommand{\CT}{\mathcal C_{3,T}}
\newcommand{\IT}{\mathfrak I_T}
\newcommand{\dd}{\,\mathrm d}
\title{Work Package B: the smoothed triple identity\\and its uniform error budget}
\author{Project proof draft}
\date{October 1, 2026}
\begin{document}
\maketitle

\noindent Task: proof and bookkeeping. Assumptions: \texttt{UNCONDITIONAL}.
Ledger status: \texttt{proved-draft}. This is the first exact identity for
Work Package B; the support-restricted asymptotic remains a target.

\section{Definitions and imported inputs}

Let \(\mathcal I\) index all occurrences of nontrivial zeros, with
multiplicity and both signs of the ordinate. Write
\(\rho_j=1/2+\delta_j+\ii\gamma_j\), so \(|\delta_j|<1/2\).
Equality of indices, complex zeros, and ordinates are three different
conditions. All powers of positive real numbers use the real logarithm.
For \(x\geq1\) and \(t\in\mathbb R\), put
\begin{align}
 D_j(t)&=1+((t-\gamma_j)+\ii\delta_j)^2,\notag\\
 \ZS(x,t)&=\sum_{j\in\mathcal I}
       \frac{2x^{\delta_j+\ii(\gamma_j-t)}}{D_j(t)},
       \label{eq:triple-S}\\
 a_x(n)&=\frac{\Lambda(n)}{\sqrt n}\min\{n/x,x/n\},\qquad
 \PS(x,t)=\sum_{n\geq1}a_x(n)e^{-\ii t\log n}.
       \label{eq:triple-P}
\end{align}
Here \(\Lambda(1)=0\). The three direct dependencies, all reconstructed
in \texttt{proofs/pair\_baseline.tex}, are:
\begin{itemize}
 \item \texttt{PAIR-EF-CONVERGENCE-001} (label \texttt{lem:ef-convergence}):
 absolute and locally uniform convergence of both series;
 \item \texttt{PAIR-COUNT-KERNEL-001} (label \texttt{lem:count-kernel}):
 \[
 Q(t):=\sum_j(1+(t-\gamma_j)^2)^{-1}\leq C_0\log(|t|+2),
 \]
 with an effective absolute constant \(C_0\);
 \item \texttt{PAIR-EF-001} (label \texttt{lem:pair-ef}): the exact identity
 \begin{equation}\label{eq:triple-EF}
 \ZS(x,t)=-\PS(x,t)+A_x(t)+H_x(t),\qquad
 A_x(t)=x^{-1}\log(|t|+2).
 \end{equation}
\end{itemize}
For precision, \(H_x\) is the following exact function, not an unspecified
error symbol. With \(s=1/2+\ii t\) and
\(\chi(w)=2^w\pi^{w-1}\sin(\pi w/2)\Gamma(1-w)\), set
\begin{align}
 H_x&=E_{\rm arch}(x,\cdot)+E_{\rm DS}(x,\cdot)
             +E_{\rm pole}(x,\cdot)+E_{\rm trivial}(x,\cdot),\notag\\
 E_{\rm arch}(x,t)&=x^{-1}\left[-\frac{\chi'}{\chi}(-1/2+\ii t)
                                     -\log(|t|+2)\right],\notag\\
 E_{\rm DS}(x,t)&=x^{-1}\frac{\zeta'}{\zeta}(3/2-\ii t),\notag\\
 E_{\rm pole}(x,t)&=\frac{2x^{1/2-\ii t}}{(1/2+\ii t)(3/2-\ii t)},\notag\\
 E_{\rm trivial}(x,t)&=2\sum_{m\geq1}
       \frac{x^{-2m-s}}{(2m+s-1)(2m+s+1)}.\label{eq:triple-H}
\end{align}
The last series is absolutely locally uniformly convergent, including
\(x=1\). The imported separate bounds, uniform in \(x\geq1,t\in\mathbb R\), are
\begin{align}
 |E_{\rm arch}|&\leq C_{\rm arch}/x,&
 |E_{\rm DS}|&\leq C_{\rm DS}/x,\notag\\
 |E_{\rm pole}|&\leq4\sqrt{x}/(1+t^2),&
 |E_{\rm trivial}|&\leq12x^{-5/2}/(|t|+2).\label{eq:triple-H-bounds}
\end{align}
All constants here are effective. Their propagation through products
will be a separate error-budget step.

\section{The exact identity}

\begin{lemma}[\texttt{TRIPLE-MASTER-001}]\label{lem:triple-master}
Fix a real nonnegative \(\omega\in C_c^\infty((1,2))\) with
\(\int_{\mathbb R}\omega(u)\dd u=1\). For \(T\geq3\), \(X,Y\geq1\),
set \(Z=XY\) and
\begin{align}
 \IT(f,g,h)&=\frac1T\int_{\mathbb R}\omega(t/T)f(t)g(t)\overline{h(t)}\dd t,
       \label{eq:triple-I}\\
 \CT(X,Y)&=\IT(\ZS(X,\cdot),\ZS(Y,\cdot),\ZS(Z,\cdot)),\notag\\
 K_T(j,k,\ell)&=\frac1T\int_{\mathbb R}
       \frac{\omega(t/T)}{D_j(t)D_k(t)\overline{D_\ell(t)}}\dd t.
       \label{eq:triple-K}
\end{align}
Then the following absolutely convergent expansion holds:
\begin{equation}\label{eq:triple-zero}
 \CT(X,Y)=8\sum_{j,k,\ell\in\mathcal I}
 X^{\rho_j+\overline{\rho_\ell}-1}
 Y^{\rho_k+\overline{\rho_\ell}-1}K_T(j,k,\ell).
\end{equation}
Its five disjoint index subseries are listed in Table~\ref{tab:diagonals}.
In particular \(\CT(X,Y)=\CT(Y,X)\), and
\begin{equation}\label{eq:triple-majorant}
 |\CT(X,Y)|\leq(8/3)^3 C_0^3XY\log^3(2T+2)\|\omega\|_1.
\end{equation}
Write \(P_x=\PS(x,\cdot)\). For a word \(w\in\{P,A,H\}^3\),
let \(N_P(w)\) be its number of letters \(P\). There is also the exact expansion
\begin{equation}\label{eq:triple-27}
 \CT(X,Y)=\sum_{w\in\{P,A,H\}^3}(-1)^{N_P(w)}
 \IT(w_1{}_X,w_2{}_Y,w_3{}_Z).
\end{equation}
Table~\ref{tab:words} records all 27 ordered terms. In the Fourier convention
\(\widehat\omega(\xi)=\int\omega(u)e^{-2\pi\ii u\xi}\dd u\),
the unsigned cubic prime term is
\begin{equation}\label{eq:triple-prime}
 \IT(P_X,P_Y,P_Z)=\sum_{m,n,k\geq1}a_X(m)a_Y(n)a_Z(k)
 \widehat\omega\!\left(\frac{T}{2\pi}\log\frac{mn}{k}\right).
\end{equation}
The series is absolutely convergent. Its exact resonance \(mn=k\) equals
\begin{equation}\label{eq:triple-resonance}
 R_{\rm res}(X,Y)=\sum_{p\ {\rm prime}}\ \sum_{a,b\geq1}
       a_X(p^a)a_Y(p^b)a_Z(p^{a+b}).
\end{equation}
If \(R_{\rm off}(X,Y,T)\) denotes the subseries in
\eqref{eq:triple-prime} with \(mn\ne k\), its contribution together with
the resonance to \(\CT\) is \(-R_{\rm res}-R_{\rm off}\).
No estimate making \(R_{\rm off}\) negligible is asserted.
\end{lemma}

\begin{proof}
For real \(v\), \(|\delta|<1/2\),
\[
 |1+(v+\ii\delta)^2|\geq1+v^2-\delta^2
                         \geq\tfrac34(1+v^2).
\]
Consequently the sum of absolute values of the individual products
in the integrand defining \(\CT\) is at most
\[
 (8/3)^3(XYZ)^{1/2}Q(t)^3
 \leq(8/3)^3XY C_0^3\log^3(|t|+2).
\]
Multiplication by \(|\omega(t/T)|/T\) gives an integrable dominating
function supported inside \((T,2T)\). Tonelli applied to absolute values
and then Fubini justify the three zero sums and their integration.
They give \eqref{eq:triple-majorant} and absolute convergence even before
integrating the absolute values of individual terms.
The numerator of a product is
\[
 8X^{\delta_j+\ii(\gamma_j-t)}
   Y^{\delta_k+\ii(\gamma_k-t)}
   Z^{\delta_\ell-\ii(\gamma_\ell-t)}
 =8X^{\delta_j+\delta_\ell+\ii(\gamma_j-\gamma_\ell)}
    Y^{\delta_k+\delta_\ell+\ii(\gamma_k-\gamma_\ell)}.
\]
This uses \(\log Z=\log X+\log Y\); all center phases cancel.
The conjugated third denominator is \(\overline{D_\ell(t)}\).
This proves \eqref{eq:triple-zero}. Exchanging the first two slots proves
the stated symmetry. Partitioning ordered index triples into their five
equality patterns proves Table~\ref{tab:diagonals}; absolute convergence
allows this regrouping without a limiting-order convention.

Substitute \eqref{eq:triple-EF} into the three slots of \(\IT\).
Each slot has three summands. The sign of a term is
\((-1)^{N_P(w)}\), including when the third slot is \(-P_Z\),
because \(-1\) is real. This proves \eqref{eq:triple-27} and
Table~\ref{tab:words}. Every function is continuous on the compact
integration window, so this finite distribution is legitimate.
Each \(H\) may in turn be expanded into its four exact named terms in
\eqref{eq:triple-H}; conjugation remains on every third-slot term.

For the prime expansion, set \(B(x)=\sum_{n\geq1}a_x(n)\).
The majorant
\(a_x(n)\leq x(\log n)n^{-3/2}\), \(n\geq2\), gives \(B(x)<\infty\).
The integral of the absolute product series is at most
\(\|\omega\|_1 B(X)B(Y)B(Z)\), which proves Fubini's hypotheses.
The oscillation is \(e^{-\ii t(\log m+\log n-\log k)}\).
Substituting \(t=Tu\) cancels \(1/T\) and gives exactly the argument
\(T\log(mn/k)/(2\pi)\) in \eqref{eq:triple-prime}.
In particular there is no extra factor of \(T\) or \(2\pi\).
Nonzero resonant coefficients require \(m,n,k\) all to be prime powers
greater than one. Unique factorization and \(mn=k\) then force
\(m=p^a,n=p^b,k=p^{a+b}\) for a single prime and \(a,b\geq1\).
Conversely every such triple is resonant. Since \(\widehat\omega(0)=1\),
\eqref{eq:triple-resonance} follows, with ordered \((a,b)\) and no
division by two. Both subseries converge absolutely by the preceding
majorant. The negative sign is the \(PPP\) row of Table~\ref{tab:words}.
\end{proof}

\section{Zero-index bookkeeping}

Set \(B_{jk\ell}=8X^{\rho_j+\overline{\rho_\ell}-1}
Y^{\rho_k+\overline{\rho_\ell}-1}K_T(j,k,\ell)\).
The last column gives counts for a finite set of \(q\) occurrences.
\begin{table}[ht]
\centering\small
\begin{tabular}{llll}
\toprule
Partition & Condition & Contribution & Count\\
\midrule
\(123\) & \(j=k=\ell\) & \(\sum_j B_{jjj}\) & \(q\)\\
\(12|3\) & \(j=k\ne\ell\) & \(\sum_{j\ne\ell}B_{jj\ell}\) & \(q(q-1)\)\\
\(13|2\) & \(j=\ell\ne k\) & \(\sum_{j\ne k}B_{jkj}\) & \(q(q-1)\)\\
\(23|1\) & \(k=\ell\ne j\) & \(\sum_{j\ne k}B_{jkk}\) & \(q(q-1)\)\\
\(1|2|3\) & all indices distinct & \(\sum_{j,k,\ell\ {\rm distinct}}B_{jk\ell}\)
 & \(q(q-1)(q-2)\)\\
\bottomrule
\end{tabular}
\caption{Disjoint index diagonals.}\label{tab:diagonals}
\end{table}

Distinct indices can represent repeated copies of the same complex zero;
distinct complex zeros can have the same ordinate. No row imposes either
of these stronger distinctions. The three partial diagonals have different
kernels and weights; they have not been replaced by three times one row.

Horizontal reflection \(\rho\mapsto1-\bar\rho\) acts by
\((\delta,\gamma)\mapsto(-\delta,\gamma)\) and preserves multiplicity.
Choose any corresponding bijection on occurrences. Simultaneous reindexing
of all three slots preserves each row and its sum, by absolute convergence.
In the reindexed summand the horizontal factor becomes
\(X^{-\delta_j-\delta_\ell}Y^{-\delta_k-\delta_\ell}\) and the denominator
inside the kernel becomes \(\overline{D_j}\,\overline{D_k}D_\ell\);
the ordinate phases are unchanged. This is an equality of full subseries,
not a termwise invariance. Reindexing only one slot preserves the unrestricted
sum but need not preserve an index diagonal. Thus retaining these horizontal
factors is compatible with functional-equation symmetry without setting
any \(\delta\) to zero.

\section{Prime, archimedean, and remainder bookkeeping}

In Table~\ref{tab:words}, each row denotes its sign times the normalized
integral of the displayed product against \(\omega(t/T)\).
All factors are evaluated at \(t\), and \(Z=XY\).
There are eight rows without \(H\) and nineteen with at least one \(H\).
The latter are retained exactly; their number is not an error estimate.

\begin{longtable}{cclc}
\caption{The 27 ordered words in \eqref{eq:triple-27}.}\label{tab:words}\\
\toprule Word & Sign & Integrand & Contains \(H\)\\\midrule
\endfirsthead
\toprule Word & Sign & Integrand & Contains \(H\)\\\midrule
\endhead
\bottomrule\endfoot
PPP & $-$ & $P_XP_Y\overline{P_Z}$ & no\\
PPA & $+$ & $P_XP_Y\overline{A_Z}$ & no\\
PPH & $+$ & $P_XP_Y\overline{H_Z}$ & yes\\
PAP & $+$ & $P_XA_Y\overline{P_Z}$ & no\\
PAA & $-$ & $P_XA_Y\overline{A_Z}$ & no\\
PAH & $-$ & $P_XA_Y\overline{H_Z}$ & yes\\
PHP & $+$ & $P_XH_Y\overline{P_Z}$ & yes\\
PHA & $-$ & $P_XH_Y\overline{A_Z}$ & yes\\
PHH & $-$ & $P_XH_Y\overline{H_Z}$ & yes\\
APP & $+$ & $A_XP_Y\overline{P_Z}$ & no\\
APA & $-$ & $A_XP_Y\overline{A_Z}$ & no\\
APH & $-$ & $A_XP_Y\overline{H_Z}$ & yes\\
AAP & $-$ & $A_XA_Y\overline{P_Z}$ & no\\
AAA & $+$ & $A_XA_Y\overline{A_Z}$ & no\\
AAH & $+$ & $A_XA_Y\overline{H_Z}$ & yes\\
AHP & $-$ & $A_XH_Y\overline{P_Z}$ & yes\\
AHA & $+$ & $A_XH_Y\overline{A_Z}$ & yes\\
AHH & $+$ & $A_XH_Y\overline{H_Z}$ & yes\\
HPP & $+$ & $H_XP_Y\overline{P_Z}$ & yes\\
HPA & $-$ & $H_XP_Y\overline{A_Z}$ & yes\\
HPH & $-$ & $H_XP_Y\overline{H_Z}$ & yes\\
HAP & $-$ & $H_XA_Y\overline{P_Z}$ & yes\\
HAA & $+$ & $H_XA_Y\overline{A_Z}$ & yes\\
HAH & $+$ & $H_XA_Y\overline{H_Z}$ & yes\\
HHP & $-$ & $H_XH_Y\overline{P_Z}$ & yes\\
HHA & $+$ & $H_XH_Y\overline{A_Z}$ & yes\\
HHH & $+$ & $H_XH_Y\overline{H_Z}$ & yes\\
\end{longtable}

\section{Endpoints, uniformity, and the next estimates}

The endpoints \(X=1\), \(Y=1\), and \(T=3\) are included. The smooth
weight vanishes in neighborhoods of \(t=T,2T\); there are no half weights.
The height window is a condition on the center \(t\), not on the zeros:
all ordinates, including \(0<|\gamma|<3\) if present, remain in the sums.
No use of \(L_T^{-1}\) or a change of local-spacing normalization is made.

Limit order is explicit: fix \(T,X,Y,\omega\); truncate each zero sum at
\(|\gamma|\leq U_r\), including endpoints with full multiplicity; then let
the three \(U_r\) tend to infinity independently. The displayed dominating
function proves convergence in any order or jointly. Prime truncations
are removed in the same way using \(B(X)B(Y)B(Z)\).
No \(T\to\infty\) limit, moving Fourier support, or shrinking smoothing
width has been interchanged with those limits. The bounds are locally
uniform in the finite parameters and explicitly retain \(\|\omega\|_1\).

Integration by parts, with no boundary terms, gives for every integer
\(N\geq1\)
\begin{equation}\label{eq:triple-decay}
 |\widehat\omega(\xi)|\leq\|\omega\|_1,\qquad
 |\widehat\omega(\xi)|\leq
 \frac{\|\omega^{(N)}\|_1}{(2\pi|\xi|)^N}\quad(\xi\ne0).
\end{equation}
For \(mn\ne k\), the second denominator becomes
\((T|\log(mn/k)|)^N\). Uniformity for a varying family of weights
requires bounds on these derivative norms. Smooth compact height support
does not supply a compact Fourier support or exact off-resonance removal.
Near \(mn=k\) a new arithmetic estimate is still needed.

As an algebraic lower-level check, replacing the second input of \(\IT\)
by the constant function \(1\) gives the smoothed pair form
\(T^{-1}\int\omega(t/T)f(t)\overline{h(t)}\dd t\).
Taking \(Y=1\) in the actual lemma does \emph{not} perform that replacement:
\(\ZS(1,t)\) remains a zero sum.

The following sections estimate the eight rows without \(H\) and propagate
the four separate bounds in \eqref{eq:triple-H-bounds} through the other
nineteen rows, before choosing any growing range for \(X,Y\). The zero
index diagonals and prime resonances are separate decompositions, with
no asserted rowwise correspondence. A Fourier-support region, an
asymptotic main term, positivity, and a new horizontal conclusion are
not established by this lemma.

\section{Uniform norm bounds and the named budget}

The task remains a proof under \texttt{UNCONDITIONAL} assumptions.
In addition to Lemma~\ref{lem:triple-master}, the direct analytic inputs are
\texttt{PAIR-EF-001}, \texttt{PAIR-PRIME-MEAN-001}, and
\texttt{PAIR-PRIME-DIAGONAL-001}, in the pair draft at labels
\texttt{lem:pair-ef}, \texttt{lem:prime-mean}, and
\texttt{lem:prime-diagonal}. The latter two give, with effective absolute
constants \(C_{\rm pair},C_D\),
\[
 \int_0^V|P_x(t)|^2\dd t\leq C_{\rm pair}V\log(2x)
 \quad(V\geq3,\ 1\leq x\leq V),\qquad
 \sum_n a_x(n)^2\leq C_D\log(2x)\quad(x\geq1).
\]
These are unconditional estimates including prime powers.

Keep the hypotheses on \(T,X,Y,Z,\omega\) from
Lemma~\ref{lem:triple-master}, and define
\begin{align}
 d\mu_T(t)&=T^{-1}\omega(t/T)\dd t,&
 q&=\log T,& L&=\log(2T+2),& \ell_x&=\log(2x),\notag\\
 W_0&=\|\omega\|_1=1,&W_\infty&=\|\omega\|_\infty,&
 W_1&=\|\omega'\|_1.&&\label{eq:uniform-parameters}
\end{align}
The measure \(\mu_T\) is a probability measure. Set
\begin{align}
 C_B&=1+\frac{2}{(1-2^{-1/2})^2},&
 C_1&=1+\frac{2}{1-2^{-1/2}},&
 C_M&=\sqrt{2C_{\rm pair}},\notag\\
 b_x&=C_B\sqrt{x}\ell_x,&\alpha_x&=L/x,&&\notag\\
 m_x&=
 \begin{cases}
 \min\{b_x,C_M\sqrt{W_\infty\ell_x}\},&1\leq x\leq2T,\\
 b_x,&x>2T,
 \end{cases}&&&&\label{eq:uniform-primitives}\\
 h_{\rm arch}(x)&=C_{\rm arch}/x,&
 h_{\rm DS}(x)&=C_{\rm DS}/x,\notag\\
 h_{\rm pole}(x)&=\frac{4\sqrt{x}}{1+T^2},&
 h_{\rm trivial}(x)&=\frac{12x^{-5/2}}{T+2},&
 h_x&=\sum_{\nu\in\mathcal J}h_\nu(x),\label{eq:uniform-remainders}
\end{align}
where \(\mathcal J=\{{\rm arch},{\rm DS},{\rm pole},{\rm trivial}\}\).
The subscripts \(x\) on \(m,h,\alpha\) suppress their displayed
dependence on \(T,\omega\).

For \(r=1,2,3\), put
\[
 J_r=\int\omega(u)\log^r(Tu+2)\dd u,\qquad
 D(a,b)=\sum_{n\geq1}a_a(n)a_b(n),\qquad
 d_{a,b}=C_D\sqrt{\ell_a\ell_b}.
\]
Integrals here and below involving \(\log(Tu+2)\) are over \((1,2)\);
the weighted functions are extended by zero outside that interval.
In particular they are smooth and compactly supported. For \(r=1,2\), set
\begin{equation}\label{eq:uniform-Omega}
 \Omega_{r,1}=\|(\omega(u)\log^r(Tu+2))'\|_1,\qquad
 V_r=W_1L^r+rW_0L^{r-1}.
\end{equation}
The notation \(D(a,b)\) refers to prime coefficients and has no relation
to the zero denominators \(D_j(t)\).

\begin{lemma}[\texttt{TRIPLE-UNIFORM-001}]\label{lem:triple-uniform}
Uniformly for \(T\geq3,\ X,Y\geq1\) and every admissible \(\omega\), the
following bounds hold, with the smoothing dependence explicitly retained:
\begin{align}
 \|P_x\|_\infty&\leq\sum_na_x(n)\leq b_x,&
 \sum_{n\geq2}\frac{a_x(n)}{\log n}&\leq C_1\sqrt{x},\notag\\
 \|P_x\|_{L^2(\mu_T)}&\leq m_x,&
 \|A_x\|_{L^\infty(T,2T)}&\leq\alpha_x,&
 \|E_\nu(x,\cdot)\|_{L^\infty(T,2T)}&\leq h_\nu(x),\notag\\
 0\leq D(a,b)&\leq d_{a,b},&
 0\leq J_r&\leq L^r,&
 \Omega_{r,1}&\leq V_r\quad(r=1,2).
 \label{eq:uniform-norms}
\end{align}
Define the retained expression
\begin{equation}\label{eq:uniform-retained}
 M_T=\frac{q^3}{Z^2}+\frac{J_1}{Y}D(X,Z)+\frac{J_1}{X}D(Y,Z)
                    -R_{\rm res}(X,Y).
\end{equation}
It is not asserted to be an asymptotic main term. There is the exact identity
\begin{align}
 \CT={}&M_T+E_{\rm archcube}
   +E_{PPA}+E_{PAA}+E_{APA}+E_{AAP}\notag\\
   &+E_{PAP,{\rm off}}+E_{APP,{\rm off}}+E_{PPP,{\rm off}}
   +\sum_{\substack{w\in\{P,A,H\}^3\\H\ {\rm occurs\ in}\ w}}E_w.
   \label{eq:uniform-budget}
\end{align}
The eight errors on the first two lines are defined and bounded in
Table~\ref{tab:uniform-eight}. In that table,
\begin{equation}\label{eq:uniform-U}
 U=\min\{b_Xm_Ym_Z,\ b_Ym_Xm_Z,\ b_Zm_Xm_Y\},\qquad
 0\leq R_{\rm res}\leq b_Xb_Yb_Z,
\end{equation}
and the complete cubic satisfies \(|\IT(P_X,P_Y,P_Z)|\leq U\).

For the remaining nineteen words, let \(x_1=X,x_2=Y,x_3=Z\), and
\(S_H(w)=\{i:w_i=H\}\). For each assignment
\(\kappa:S_H(w)\to\mathcal J\), replace the corresponding factors by
\(E_{\kappa(i)}(x_i,\cdot)\), retaining \(P,A\) elsewhere, to get
functions \(f_1,f_2,f_3\). Define
\begin{equation}\label{eq:uniform-refinement}
 E_{w;\kappa}=(-1)^{N_P(w)}\IT(f_1,f_2,f_3),\qquad
 E_w=\sum_\kappa E_{w;\kappa}.
\end{equation}
Then, separately for every assignment,
\begin{equation}\label{eq:uniform-component-bound}
 |E_{w;\kappa}|\leq
 \prod_{\{i:w_i=P\}}m_{x_i}
 \prod_{\{i:w_i=A\}}\alpha_{x_i}
 \prod_{\{i:w_i=H\}}h_{\kappa(i)}(x_i).
\end{equation}
Summing these nonnegative bounds gives
\begin{equation}\label{eq:uniform-word-bound}
 |E_w|\leq
 \prod_{\{i:w_i=P\}}m_{x_i}
 \prod_{\{i:w_i=A\}}\alpha_{x_i}
 \prod_{\{i:w_i=H\}}h_{x_i}.
\end{equation}
There are \(208\) refined remainder terms, grouped into nineteen words;
the refined terms and their aggregates are two descriptions of the same
summands. All constants in these assertions are effective and independent
of \(T,X,Y,\omega\), apart from the displayed norms.
\end{lemma}

\begin{table}[p]
\centering\small
\begin{tabular}{p{0.23\textwidth}p{0.36\textwidth}p{0.32\textwidth}}
\toprule
Name & Exact signed definition & Absolute bound\\
\midrule
\(E_{\rm archcube}\) &
 \((J_3-q^3)/Z^2\) &
 \((L^3-q^3)/Z^2\)\\[5pt]
\(E_{PPA}\) & \(\IT(P_X,P_Y,A_Z)\) &
 \(b_Xb_Y V_1/(ZT\log4)\)\\[5pt]
\(E_{PAA}\) & \(-\IT(P_X,A_Y,A_Z)\) &
 \(C_1\sqrt{X}V_2/(TYZ)\)\\[5pt]
\(E_{APA}\) & \(-\IT(A_X,P_Y,A_Z)\) &
 \(C_1\sqrt{Y}V_2/(TXZ)\)\\[5pt]
\(E_{AAP}\) & \(-\IT(A_X,A_Y,P_Z)\) &
 \(C_1\sqrt{Z}V_2/(TXY)\)\\[5pt]
\(E_{PAP,{\rm off}}\) &
 \(\IT(P_X,A_Y,P_Z)-J_1D(X,Z)/Y\) &
 \((L/Y)(m_Xm_Z+d_{X,Z})\)\\[5pt]
\(E_{APP,{\rm off}}\) &
 \(\IT(A_X,P_Y,P_Z)-J_1D(Y,Z)/X\) &
 \((L/X)(m_Ym_Z+d_{Y,Z})\)\\[5pt]
\(E_{PPP,{\rm off}}\) &
 \(-\IT(P_X,P_Y,P_Z)+R_{\rm res}\) &
 \(U+R_{\rm res}\leq U+b_Xb_Yb_Z\)\\
\bottomrule
\end{tabular}
\caption{Named errors from the eight words without \(H\).
Each \(V_r\) may be replaced by the sharper \(\Omega_{r,1}\).
The mixed-prime and cubic off-diagonal bounds make no assertion of
cancellation.}\label{tab:uniform-eight}
\end{table}

\begin{proof}
For \(n\leq x\), the inequality \(\Lambda(n)\leq\log n\) gives
\[
 \sum_{2\leq n\leq x}a_x(n)
 \leq x^{-1}\sum_{2\leq n\leq x}\sqrt n\log n
 \leq\sqrt{x}\ell_x.
\]
For \(j\geq0\), the interval \(2^jx<n\leq2^{j+1}x\) contains at most
\(2^{j+1}x\) positive integers. Consequently its contribution is at most
\[
 2\sqrt{x}\,2^{-j/2}\log(2^{j+1}x)
 \leq2\sqrt{x}\ell_x(j+1)2^{-j/2}.
\]
Summing the geometric series and its derivative gives \(b_x\) as stated.
Using \(\Lambda(n)/\log n\leq1\) in the same argument removes the
factor \((j+1)\ell_x\), giving \(C_1\sqrt{x}\).
The split counts an integer \(n=x\) exactly once and applies also to
\(1\leq x<2\). All series are nonnegative and convergent.

Since \(\mu_T\) has mass one, \(\|P_x\|_{L^2(\mu_T)}\leq b_x\).
If \(x\leq2T\), positivity of the weight and the imported mean-square bound
with \(V=2T\) give
\[
 \int|P_x|^2\dd\mu_T
 \leq \frac{W_\infty}{T}\int_0^{2T}|P_x(t)|^2\dd t
 \leq2C_{\rm pair}W_\infty\ell_x.
\]
This proves the piecewise bound, including \(x=2T\). No extension of
the imported mean-square theorem to \(x>2T\) is used.
The \(A\) and four remainder bounds follow pointwise from
\eqref{eq:triple-H-bounds} on \(T<t<2T\). Cauchy--Schwarz for the
coefficient sequences and the imported diagonal bound give
\(D(a,b)\leq\sqrt{D(a,a)D(b,b)}\leq d_{a,b}\).

On the support of \(\omega\), \(q\leq\log(Tu+2)\leq L\) and
\(T/(Tu+2)\leq1\). Thus \(q^r\leq J_r\leq L^r\), and the product rule
gives
\[
 \Omega_{r,1}\leq
 \|\omega'\|_1L^r+r\|\omega\|_1L^{r-1}=V_r.
\]
There are no endpoint derivative contributions: \(\omega\) vanishes
in neighborhoods of \(1,2\). In particular
\(0\leq E_{\rm archcube}\leq(L^3-q^3)/Z^2\).

For the oscillating terms, set
\[
 F_r(\lambda)=\int\omega(u)\log^r(Tu+2)e^{-\ii Tu\lambda}\dd u .
\]
One integration by parts gives
\begin{equation}\label{eq:uniform-ibp}
 |F_r(\lambda)|\leq\frac{\Omega_{r,1}}{T|\lambda|}
 \quad(\lambda\ne0).
\end{equation}
In project Fourier units \(F_r(\lambda)\) is the transform of
\(\omega(u)\log^r(Tu+2)\) at \(T\lambda/(2\pi)\).
The \(2\pi\) from differentiation cancels against this denominator;
the normalized substitution \(t=Tu\) introduces no further factor.
Since all prime coefficient sums are absolutely convergent and the
weights bounded, Fubini applies at fixed parameters.

The \(PPA\) integral is
\[
 \frac1Z\sum_{m,n\geq2}a_X(m)a_Y(n)F_1(\log(mn)).
\]
Here \(\log(mn)\geq\log4\), which proves its table bound.
The \(PAA,APA,AAP\) integrals have respective frequencies
\(\log n,\log n,-\log n\) and scalar denominators \(YZ,XZ,XY\).
Apply \eqref{eq:uniform-ibp} and the bound for
\(\sum a_x(n)/\log n\). The third of these terms still has its conjugation;
the change of frequency sign does not affect the absolute bound.

For example the mixed term is exactly
\[
 \IT(P_X,A_Y,P_Z)=
 \frac1Y\sum_{m,k\geq2}a_X(m)a_Z(k)F_1(\log(m/k)).
\]
The diagonal is \(J_1D(X,Z)/Y\). Cauchy--Schwarz bounds the
full integral in absolute value by \((L/Y)m_Xm_Z\).
Subtracting its diagonal gives the bound in the table by the triangle
inequality. The \(APP\) calculation is identical with \(X,Y\) interchanged.
These bounds do not treat the near frequencies as separated from zero.

For the cubic integral take any one factor in sup norm and the other two
in \(L^2(\mu_T)\). The three choices give \(U\).
The resonant sum is a nonnegative subseries of the product of the three
coefficient sums, so \(R_{\rm res}\leq b_Xb_Yb_Z\). Its signed
off-resonance error is \(-R_{\rm off}\) from Lemma~\ref{lem:triple-master};
the triangle inequality gives \(U+R_{\rm res}\).

Each word containing \(H\) has at most two \(P\) slots. Bound all other
factors in sup norm, and use Cauchy--Schwarz for two prime factors or
\(\|P_x\|_{L^1(\mu_T)}\leq\|P_x\|_{L^2(\mu_T)}\) for one.
This proves \eqref{eq:uniform-component-bound}. Finite distributivity
proves \eqref{eq:uniform-word-bound}, with no absorbed constants.
There are twelve words with one \(H\), six with two, and one with three:
\[
 12\cdot4+6\cdot4^2+4^3=208.
\]
Finally substitute the exact signed definitions into
\eqref{eq:triple-27}. The \(AAA\) term is \(J_3/(XYZ)=J_3/Z^2\).
Extracting its \(q^3/Z^2\), the two mixed diagonals, and the negatively
signed cubic resonance proves \eqref{eq:uniform-budget}.
\end{proof}

\section{Range, interpretation, and remaining losses}

All estimates in Lemma~\ref{lem:triple-uniform} hold on the entire range
\(T\geq3,\ X,Y\geq1\), using the specified branch of \(m_x\).
If \(XY\leq2T\), all three prime factors admit the sharper branch since
\(1\leq X,Y\leq XY\); equality \(XY=2T\) is included.
The line \(x=2T\) is a boundary of an estimate's applicability, not a
Fourier-support boundary. For \(x>2T\) the elementary bound remains valid.
The norm majorant need not be continuous when its branch changes.

Every zero occurrence from Lemma~\ref{lem:triple-master} is still retained.
Its symmetry under simultaneous horizontal reflection therefore persists;
estimating its exactly equal prime-side expression does not set
\(\delta_\rho=0\). Swapping \(X,Y\) interchanges \(PAA,APA\) and the two
mixed off-diagonal errors; the other definitions transform by their
ordered slots. Setting a slot function to \(1\) recovers the ordinary
weighted pair Cauchy--Schwarz bound. Setting its base to \(1\) does not
remove that slot.

The limit order remains fixed \(T,X,Y,\omega\), followed by removal of
absolutely convergent prime or zero cutoffs. Bounds involving \(\omega'\)
are used only after the weighted integrals are defined; no differentiation
of an infinite prime series is taken. All displayed estimates are then
uniform on their stated ranges. No shrinking-width limit is taken.
Varying weights must retain \(W_\infty,W_1\); mass normalization alone
does not bound them.

The pure archimedean term and the one-prime and same-sign two-prime terms
have elementary bounds, including the displayed \(1/T\) oscillatory gain.
The mixed-prime diagonal bound imports PNT-level information through
\texttt{PAIR-PRIME-DIAGONAL-001}. The weighted prime norms import the
unconditional mean-value and prime-counting inputs in
\texttt{PAIR-PRIME-MEAN-001}. No new prime-pair or prime-triple conjecture is
used. The mixed-prime and cubic off-diagonal bounds are triangle-inequality
bounds and do not yield a uniform negligible error. No rowwise matching
with zero-index diagonals has been established.

The file \path{research/triple-error-budget.md} records the nineteen
remainder-containing rows and their component assignments. They form one budget:
an aggregate error and its component terms must never both be added.
There is no height-truncation, desmoothing, or support-boundary error in
this identity. The next analytic task is to sharpen the mixed-prime and
cubic off-diagonal bounds and assess the remainder scales on a chosen
growing parameter range; the present bounds alone do not settle an
asymptotic or supply horizontal information.
\end{document}
